\chapter{Related Work}
\label{ch:related}

The detection of malicious packages in open-source registries has attracted considerable research attention in recent years, 
and the proposed approaches can be broadly grouped into four categories: 
static analysis, dynamic analysis, machine learning and LLM-based methods, and metric-based analysis.
\cref{sec:staticAnalysis} examines representative static analysis tools, 
while \cref{sec:dynamicAnalysis} covers dynamic analysis approaches. 
\cref{sec:MLAndLLM} describes in detail the two most representative LLM-based
tools, \texttt{SpiderScan} (\cref{subsec:SpiderScan}) and \texttt{SocketAI} (\cref{subsec:SocketAI}), 
and \cref{sec:metricBasedAnalysis} surveys metric-based methods that rely on measurable structural properties of packages.
Finally, \cref{sec:positioning} discusses the two main limitations affecting most existing tools which motivate this research, 
and clarifies how this thesis relates to and differs from existing works, highlighting the contribution it makes to the field.

\section{Static Analysis}
\label{sec:staticAnalysis}

Static analysis examines the code and metadata of a package without executing it. 
Its main advantage is speed: packages can be screened in seconds without requiring
a sandboxed execution environment. 
Several tools in this space rely on a predefined set of sensitive \acp{API}
to flag suspicious packages, inspecting whether dangerous functions
are called or how data flows between sources and sinks.

One representative of this family is \texttt{CodeQL}, a query-based static analysis engine.
It works by extracting the static features of the code into a searchable representation,
treating the codebase as a database that can be interrogated using specific queries.
\textcite{zahan2025leveraginglargelanguagemodels} used thirty-nine custom \texttt{CodeQL} 
rules, originally developed by Froh et al.~\cite{unknown}, to target malicious JavaScript behaviour.
These rules monitor how code interacts with sensitive system \acp{API} by tracking data flows 
from untrusted sources to dangerous sinks, covering patterns such as credential exfiltration through network \acp{API}, 
connections to suspicious domains, unauthorized file system access, and the execution of dynamically generated code via sinks like \texttt{eval()}. 
Their work is presented as part of the \texttt{SocketAI} workflow described further in \cref{subsec:SocketAI}, and it 
establishes \texttt{CodeQL}-based static analysis as a practical baseline for comparison with more advanced techniques. 

A fundamental shortcoming of purely static approaches, however, is their difficulty
in handling obfuscated code, as described in \cref{subsec:ObfuscationTechniques}.
When dangerous behavior is deliberately hidden behind encoding or renaming techniques, 
it cannot be matched against a fixed set of syntactic rules, 
causing static analyzers to miss the threat entirely.

\section{Dynamic Analysis}
\label{sec:dynamicAnalysis}

Dynamic analysis overcomes the obfuscation problem by executing the package in a controlled environment 
and observing its runtime behaviour. 
The primary cost is the significant overhead of setting up and tearing down sandbox environments for every candidate package, 
as well as the risk that some malicious payloads may not be triggered during the bounded execution window.

\texttt{DONAPI}~\cite{huang2024donapimaliciousnpmpackages} is a representative tool 
that combines static and dynamic analysis in a two-phase pipeline.
In the static phase, the entry-point files of a package and their dependencies are analyzed 
to extract behavioral features based on \ac{API} call sequences.
Rather than checking for isolated dangerous functions, \texttt{DONAPI} focuses on the order in which \acp{API} are invoked, 
capturing combinations that are characteristic of malicious activity.
The authors identified twelve specific indicators that describe typical malicious behaviors, 
such as sending sensitive information to an external server or downloading and executing content as a string.
These features are used to assign a suspiciousness score and to filter which packages proceed to the more expensive dynamic phase.
Packages flagged as suspicious are then executed inside isolated Docker containers, 
where a custom instrumentation layer intercepts and records runtime \ac{API} calls. 
Rather than relying on system-call tracing tools such as \texttt{Strace} or \texttt{Sysdig}, 
which operate at the OS level and lack interpretability, 
\texttt{DONAPI} injects monitoring code directly into 132 Node.js \ac{API} implementations,
capturing both the calls and their explicit parameters without interference from unrelated processes running on the same host.

After deployment in a real-world monitoring environment from January to June 2023, 
\texttt{DONAPI} identified and manually confirmed 325 previously unknown malicious packages, 
and discovered 246 \ac{API} call sequences not seen in prior samples. 
Despite these results, the authors of \texttt{SpiderScan}~\cite{10.1145/3691620.3695492} 
note that the heavyweight nature of the dynamic execution phase, 
requiring a full container per package, significantly impedes \texttt{DONAPI}'s throughput and practical scalability.

\section{Machine Learning and LLM-based Methods}
\label{sec:MLAndLLM}

The third category encompasses techniques that use machine learning models, including \acp{LLM}, 
to detect malicious packages and identify malicious code fragments within them.
The two most representative tools in this space are \texttt{SpiderScan} and \texttt{SocketAI}, 
described in the following subsections.

\subsection{SpiderScan}
\label{subsec:SpiderScan}

\texttt{SpiderScan}~\cite{10.1145/3691620.3695492} proposes a graph-based approach that avoids the cost of full dynamic 
execution while capturing richer semantic information than plain static analysis. 
The tool operates in two phases: an offline phase to construct malicious behavior models from known samples, 
and an online phase to screen newly published packages against those models.

In the offline phase, \texttt{SpiderScan} processes a set of known malicious packages and constructs a 
\emph{malicious behavior graph (MBG)} for each distinct attack pattern. 
Each graph node represents a sensitive \ac{API} call, and edges encode two types of relationships: 
\emph{control-flow dependencies}, capturing the conditions under which a call is made, 
and \emph{data-flow dependencies}, which capture the variables that carry data between calls.
Rather than limiting the sensitive \acp{API} to a hand-curated fixed list, 
\texttt{SpiderScan} uses an \ac{LLM} to recognize sensitive calls in both Node.js built-in modules and third-party dependencies, 
which are too numerous and too rapidly evolving to enumerate manually.

In the online phase, given a target package, \texttt{SpiderScan} constructs its
\emph{suspicious behavior graphs (SBGs)} and matches them against the \emph{MBGs}.
Multiple \emph{SBGs} are generated because a package may contain several distinct
behavioral patterns, each requiring separate graph representation.
Only when a match is found does \texttt{SpiderScan} invoke dynamic analysis and an additional \ac{LLM} query to confirm maliciousness,
keeping the expensive components on the critical path only when strictly necessary.

In its evaluation, \texttt{SpiderScan} monitored 298,504 newly published packages in \ac{NPM} over three months,
identifying 249 previously unknown malicious cases.
A limitation noted by its authors is that \ac{LLM} queries for \ac{API} recognition 
and maliciousness confirmation carry non-trivial costs per query, which become significant at registry-level monitoring scale.

\subsection{SocketAI}
\label{subsec:SocketAI}

\textcite{zahan2025leveraginglargelanguagemodels} propose \texttt{SocketAI}, 
a malicious code review workflow centred on \texttt{GPT-3} and \texttt{GPT-4}. 
The design relies on a two-part prompt structure:
a \emph{system-role} prompt sets the behavioral context for the model,
instructing it to act as a security expert reviewing JavaScript code for malicious patterns, 
while a \emph{user-role} prompt provides the file content to be analyzed. 
To contain \ac{API} costs, the workflow uses \texttt{CodeQL}-based static analysis as a
pre-screening step that reduces the number of files forwarded to the \ac{LLM}.

Evaluated on 5,115 \ac{NPM} packages, of which 2,180 contain malicious code, 
\texttt{GPT-4} outperforms \texttt{CodeQL} alone, 
while the combined static-\ac{LLM} pipeline achieves the best cost-accuracy trade-off overall.
As with \texttt{SpiderScan}, the use of commercial \ac{LLM} \acp{API} makes continuous monitoring of the entire registry expensive,
representing a shared limitation of both approaches.

\section{Metric-based Analysis}
\label{sec:metricBasedAnalysis}

A complementary line of research focuses on measurable, structural characteristics of 
source code files and package artifacts, rather than on the semantics of \ac{API} calls.
This approach does not require execution or language model inference,
making it lightweight and straightforward to apply at scale.

\textcite{weissberg2025llm} study the relationship between code metrics and vulnerable code, 
proposing a taxonomy organized into several families: 
\emph{code smell indicators} such as the number of \texttt{goto} statements and function pointers, 
\emph{complexity metrics} such as the number of loops and local variables,
\emph{memory management operation counts}, and \emph{Abstract Syntax Tree (AST) metrics}. 
While their focus is on vulnerability identification rather than \ac{SSCA} detection, 
their metric taxonomy provides a broad reference point for what structural properties can be extracted from source code.

\textcite{zoratti2025valutazione} addresses a narrower but directly supply-chain-relevant question: 
can simple metrics derived from the whitespace and character distribution 
of source files detect \ac{BST} and \ac{HUT} attacks?
\ac{BST} attacks hide malicious code after long sequences of blank characters, exploiting parser behaviors,
while \ac{HUT} attacks use invisible Unicode control characters, such as right-to-left override markers,
to make code appear different to a human reviewer than what is actually parsed and executed. 
\textcite{zoratti2025valutazione} proposed the ratio of whitespace to printable 
characters as a key metric to detect such anomalies in GitHub repositories.
A significant aspect of his methodology is that the analysis is performed across multiple consecutive versions of a project,
studying how the metric evolves over time rather than examining a single snapshot. 
This multi-version perspective is shared with the present thesis, 
which similarly analyzes the evolution of metrics across 20 consecutive package versions to identify anomalous releases.

\section{Positioning of This Work}
\label{sec:positioning}

Despite the breadth of the existing approaches, 
two recurring limitations affect most of the tools discussed above.

The first is the \emph{absence of a longitudinal perspective}.
Most tools, including \texttt{DONAPI}~\cite{huang2024donapimaliciousnpmpackages} and \texttt{SpiderScan}~\cite{10.1145/3691620.3695492},
analyze only the most recently published version of a package, without tracking how its properties change across releases. 
This design choice makes it impossible to detect an abrupt anomalous change in 
a metric that is otherwise stable, which is precisely the signature left by the 
\acp{SSCA} analyzed in \cref{ch:AnalysisRecentAttacks}.

The second is \emph{operational cost}. 
Dynamic analysis, as in \texttt{DONAPI}, requires container instantiation per candidate package version,
while \ac{LLM}-based methods, as in \texttt{SpiderScan} and \texttt{SocketAI}~\cite{zahan2025leveraginglargelanguagemodels}, 
incur per-query inference costs that grow with monitoring volume and introduce dependencies on external services.
Together, these constraints make continuous, registry-scale monitoring challenging.

The present thesis occupies a complementary position in this landscape.
Rather than conducting semantic analysis of \ac{API} calls 
or relying on dynamic execution or commercial \ac{LLM} services, 
it evaluates a set of lightweight, easily computable metrics derived directly from tarball artifacts:
\emph{file types}, \emph{package size}, \emph{longest line length}, \emph{unique hexadecimal values} and \emph{unique Ethereum addresses}.
Their effectiveness is evaluated against a baseline established from 629 popular goodware packages, 
and their evolution is tracked across 20 consecutive versions per package.
As shown in \cref{ch:AnalysisRecentAttacks}, the four recent \acp{SSCA} analyzed in this work 
leave measurable traces in at least one of these metrics, 
suggesting that a metric-based, multi-version monitoring strategy can serve as a
practical and low-cost complement to the more semantically rich approaches surveyed above.